% Source: about_the_equivalence_between_CE_and_MLE.tex
\begin{notechapter}
\chapter{Cross-Entropy Minimization and Maximum Likelihood}\label{ch:ce-mle-equivalence}

\begin{abstract}
    This document provides a detailed explanation of the fundamental equivalence between cross-entropy minimization and maximum likelihood estimation. We prove this equivalence mathematically, provide intuitive explanations, demonstrate it with concrete examples, and explore its implications across various machine learning scenarios including classification, regression, and other probabilistic models.
\end{abstract}

\section{Introduction}

In machine learning, particularly in classification tasks, the cross-entropy loss function is ubiquitous. Similarly, in statistics, maximum likelihood estimation is a cornerstone principle for parameter estimation. This document explores the profound connection between these two concepts and why this equivalence makes cross-entropy such a powerful loss function.

\section{Mathematical Foundation}

\subsection{Maximum Likelihood Estimation (MLE)}

Given a dataset $D = \{(x_1, y_1), (x_2, y_2), \ldots, (x_n, y_n)\}$ and a parametric model that outputs a probability distribution $Q_\theta(Y|X)$, the likelihood function is:

\begin{definition}[Likelihood Function]
\[
\mathcal{L}(\theta) = \prod_{i=1}^n Q_\theta(y_i | x_i)
\]
\end{definition}

The maximum likelihood estimator seeks parameters that make the observed data most probable:

\begin{definition}[Maximum Likelihood Estimator]
\[
\hat{\theta}_{MLE} = 
\operatorname*{arg\,max}_{\theta} \mathcal{L}(\theta) = 
\operatorname*{arg\,max}_{\theta} \sum_{i=1}^n \log Q_\theta(y_i | x_i)
\]
\end{definition}

\subsection{Empirical Cross-Entropy}

The cross-entropy between the true data distribution $P$ and our model $Q_\theta$ is:

\begin{definition}[Empirical Cross-Entropy]
\[
\hat{H}(P, Q_\theta) = -\frac{1}{n} \sum_{i=1}^n \log Q_\theta(y_i | x_i)
\]
\end{definition}

The cross-entropy minimizer is:

\begin{definition}[Cross-Entropy Minimizer]
\[
\hat{\theta}_{CE} = 
\operatorname*{arg\,min}_{\theta} \hat{H}(P, Q_\theta) = 
\operatorname*{arg\,min}_{\theta} \left[ -\frac{1}{n} \sum_{i=1}^n \log Q_\theta(y_i | x_i) \right]
\]
\end{definition}

\section{The Fundamental Equivalence}

\begin{theorem}[Equivalence of MLE and Cross-Entropy Minimization]
For any probabilistic model $Q_\theta(Y|X)$ and dataset $D$, the maximum likelihood estimator is identical to the cross-entropy minimizer:
\[
\hat{\theta}_{MLE} = \hat{\theta}_{CE}
\]
\end{theorem}

\begin{proof}
\begin{align*}
\hat{\theta}_{CE} &= \operatorname*{arg\,min}_{\theta} \left[ -\frac{1}{n} \sum_{i=1}^n \log Q_\theta(y_i | x_i) \right] \\
&= \operatorname*{arg\,max}_{\theta} \left[ \sum_{i=1}^n \log Q_\theta(y_i | x_i) \right] \quad \text{(multiplying by -1 flips argmin to argmax)} \\
&= \operatorname*{arg\,max}_{\theta} \left[ \prod_{i=1}^n Q_\theta(y_i | x_i) \right] \quad \text{(log is monotonic)} \\
&= \hat{\theta}_{MLE}
\end{align*}

The factor $\frac{1}{n}$ does not affect the optimization since it's a positive constant.
\end{proof}

\section{Intuitive Interpretation}

\subsection{The "Goodness of Fit" Perspective}

\textbf{Maximum Likelihood}: "Find the parameters that make the observed data most probable."

\textbf{Cross-Entropy Minimization}: "Find the parameters that minimize the surprise when we see the true data."

These are two sides of the same coin:
\begin{itemize}
    \item High probability assigned to true labels → High likelihood AND Low cross-entropy
    \item Low probability assigned to true labels → Low likelihood AND High cross-entropy
\end{itemize}

\subsection{The "Penalty" Interpretation}

The term $-\log Q_\theta(y_i | x_i)$ acts as a penalty:
\begin{itemize}
    \item If $Q_\theta(y_i | x_i) \to 1$: $-\log Q_\theta(\cdot) \to 0$ (no penalty)
    \item If $Q_\theta(y_i | x_i) \to 0$: $-\log Q_\theta(\cdot) \to \infty$ (infinite penalty)
    \item The penalty grows steeply as predictions become more incorrect
\end{itemize}

\section{Concrete Example: Binary Classification}

\subsection{Scenario}
Consider 3 data points: $(x_1, y_1=1), (x_2, y_2=0), (x_3, y_3=1)$ with model predictions: $p_1 = 0.8, p_2 = 0.3, p_3 = 0.9$.

\subsection{Calculations}

\textbf{Log-Likelihood}:
\begin{align*}
\mathcal{LL} &= \log(0.8) + \log(1-0.3) + \log(0.9) \\
&= \log(0.8) + \log(0.7) + \log(0.9) \\
&\approx -0.223 - 0.357 - 0.105 = -0.685
\end{align*}

\textbf{Cross-Entropy}:
\[
\hat{H} = -\frac{1}{3}[\log(0.8) + \log(0.7) + \log(0.9)] = -\frac{1}{3}(-0.685) \approx 0.228
\]

\subsection{Optimization Direction}
\begin{itemize}
    \item To \textbf{maximize} log-likelihood: push $-0.685$ closer to $0$
    \item To \textbf{minimize} cross-entropy: push $0.228$ closer to $0$
    \item Both objectives push $p_1, p_2, p_3$ toward the true labels $(1, 0, 1)$
\end{itemize}

\section{Theoretical Implications}

\subsection{Statistical Properties}

By equivalencing cross-entropy minimization with MLE, we inherit all desirable properties of maximum likelihood estimation:

\begin{itemize}
    \item \textbf{Consistency}: As $n \to \infty$, $\hat{\theta} \to \theta^*$ (true parameters)
    \item \textbf{Asymptotic Normality}: $\sqrt{n}(\hat{\theta} - \theta^*) \xrightarrow{d} \mathcal{N}(0, I^{-1}(\theta^*))$
    \item \textbf{Efficiency}: Achieves the Cramer-Rao lower bound (minimum variance)
    \item \textbf{Invariance}: If $\hat{\theta}$ is MLE for $\theta$, then $g(\hat{\theta})$ is MLE for $g(\theta)$
\end{itemize}

\subsection{Information-Theoretic Connection}

When the true distribution $P$ is the empirical distribution of our data:
\[
\hat{H}(P, Q_\theta) = D_{KL}(P \parallel Q_\theta) + H(P)
\]
where $H(P)$ is constant. Therefore:
\[
\operatorname*{arg\,min}_{\theta} \hat{H}(P, Q_\theta) = \operatorname*{arg\,min}_{\theta} D_{KL}(P \parallel Q_\theta)
\]

This shows that cross-entropy minimization is equivalent to minimizing KL divergence to the empirical distribution.

\section{Special Cases and Extensions}

\subsection{Linear Regression with Gaussian Noise}

Assume $y_i \sim \mathcal{N}(f_\theta(x_i), \sigma^2)$ where $f_\theta(x_i) = \theta^T x_i$.

The log-likelihood is:
\begin{align*}
\log \mathcal{L}(\theta) &= \sum_{i=1}^n \log \left[ \frac{1}{\sqrt{2\pi\sigma^2}} \exp\left(-\frac{(y_i - f_\theta(x_i))^2}{2\sigma^2}\right) \right] \\
&= -\frac{n}{2} \log(2\pi\sigma^2) - \frac{1}{2\sigma^2} \sum_{i=1}^n (y_i - f_\theta(x_i))^2
\end{align*}

Maximizing log-likelihood is equivalent to minimizing Mean Squared Error (MSE):
\[
\operatorname*{arg\,max}_{\theta} \log \mathcal{L}(\theta) = \operatorname*{arg\,min}_{\theta} \frac{1}{n} \sum_{i=1}^n (y_i - f_\theta(x_i))^2
\]

\subsection{Poisson Regression}

For count data where $y_i \sim \text{Poisson}(\lambda_i)$ with $\lambda_i = \exp(\theta^T x_i)$:

The probability mass function is:
\[
P(y_i|\lambda_i) = \frac{\lambda_i^{y_i} e^{-\lambda_i}}{y_i!}
\]

The log-likelihood is:
\begin{align*}
\log \mathcal{L}(\theta) &= \sum_{i=1}^n \left[ y_i \log(\lambda_i) - \lambda_i - \log(y_i!) \right] \\
&= \sum_{i=1}^n \left[ y_i (\theta^T x_i) - \exp(\theta^T x_i) - \log(y_i!) \right]
\end{align*}

Maximizing this likelihood leads to a different loss function than cross-entropy or MSE, tailored to count data. The corresponding loss function (negative log-likelihood) is:
\[
L(\theta) = \sum_{i=1}^n \left[ \exp(\theta^T x_i) - y_i (\theta^T x_i) + \log(y_i!) \right]
\]

Note that $\log(y_i!)$ is constant with respect to $\theta$ and can be omitted during optimization.
\subsection{Logistic Regression}

For binary classification with $y_i \in \{0,1\}$ and $p_i = \sigma(\theta^T x_i) = \frac{1}{1 + \exp(-\theta^T x_i)}$:

The log-likelihood is:
\[
\log \mathcal{L}(\theta) = \sum_{i=1}^n \left[ y_i \log(p_i) + (1-y_i) \log(1-p_i) \right]
\]

This is exactly the binary cross-entropy loss (negated).

\subsection{Multinomial Regression}

For multi-class classification with $K$ classes, using softmax:
\[
p_{i,k} = \frac{\exp(\theta_k^T x_i)}{\sum_{j=1}^K \exp(\theta_j^T x_i)}
\]

The log-likelihood is:
\[
\log \mathcal{L}(\theta) = \sum_{i=1}^n \sum_{k=1}^K \mathbb{I}(y_i = k) \log(p_{i,k})
\]

This is exactly the categorical cross-entropy loss.

\subsection{Bernoulli and Categorical Distributions}

The cross-entropy loss naturally emerges from assuming:
\begin{itemize}
    \item \textbf{Binary case}: $y_i \sim \text{Bernoulli}(p_i)$
    \item \textbf{Multi-class case}: $y_i \sim \text{Categorical}(p_{i,1}, \ldots, p_{i,K})$
\end{itemize}

\subsection{Exponential Family Distributions}

For any distribution in the exponential family:
\[
p(y|\eta) = h(y) \exp(\eta^T T(y) - A(\eta))
\]

The negative log-likelihood is:
\[
-\log \mathcal{L}(\theta) = \sum_{i=1}^n \left[ A(\eta_i) - \eta_i^T T(y_i) \right] + \text{constant}
\]

This provides a generalized "cross-entropy" for any exponential family distribution.

\subsection{Robust Regression with Laplace Noise}

If we assume $y_i \sim \text{Laplace}(f_\theta(x_i), b)$ instead of Gaussian:

The log-likelihood becomes:
\[
\log \mathcal{L}(\theta) = -n \log(2b) - \frac{1}{b} \sum_{i=1}^n |y_i - f_\theta(x_i)|
\]

Maximizing this is equivalent to minimizing Mean Absolute Error (MAE), which is more robust to outliers.

\section{Practical Implications}

\subsection{Why This Matters}

The equivalence explains why cross-entropy is so effective:

\begin{itemize}
    \item \textbf{Not arbitrary}: It's the unique loss (up to constants) corresponding to maximum likelihood
    \item \textbf{Theoretically grounded}: Inherits optimality properties from MLE
    \item \textbf{Well-calibrated}: Produces proper probability estimates
    \item \textbf{Universally applicable}: Works for any probabilistic model
\end{itemize}

\subsection{Limitations and Considerations}

\begin{itemize}
    \item MLE can overfit with small datasets
    \item May be sensitive to outliers (depending on the distributional assumption)
    \item Requires the model to be well-specified
    \item In high dimensions, regularization (MAP estimation) is often preferred
\end{itemize}

\section{Conclusion}

The equivalence $\text{Minimizing cross-entropy} \equiv \text{Maximizing log-likelihood}$ is profound because:

\begin{itemize}
    \item It connects optimization to fundamental statistical principles
    \item It explains the theoretical optimality of cross-entropy
    \item It provides a unified framework for understanding different loss functions
    \item It reveals that many common loss functions (MSE, MAE, cross-entropy) emerge naturally from specific distributional assumptions
\end{itemize}

This equivalence is why cross-entropy isn't just a "good" loss function—it's the \textbf{canonical} loss function for probabilistic classification and beyond.
\end{notechapter}
